% Standalone section fragment. Requires amsmath, amssymb, amsthm, hyperref.
% Ordinary mathematical proof; deliberately does not claim a new Lean build.
\section{Why the consistency equivalence is provable in PA}
\label{sec:pa-bridge-audit}

This section records an ordinary mathematical proof, with its interfaces to
the literal machine made explicit.  A new end-to-end Lean or Rocq certificate
is not asserted.  The distinction matters: absence of a packaged proof-assistant
declaration is not, by itself, an obstruction to an ordinary proof that a
sentence is derivable in PA.  Conversely, testing a long blank-tape prefix
cannot replace the uniform simulation hypotheses below.

The proof-producing binary-comprehension argument and the diagonal-vacuity
calculus are prior work in the supplied archive collection.  In particular,
the reports \texttt{pa\_tm\_research\_2026-11-13} and
\texttt{pa\_consistency\_reduction} contain the direct arithmetic translation;
\texttt{pa\_consistency\_comparison\_and\_new196\_2026-11-13} contains the
diagonal and separate-head constructions.  The present account checks their
interfaces against the pinned source and describes a useful simplification:
the reverse invariant can carry PA proofs directly, so it does not need a
new identification of the auxiliary native checker's numerical certificate
encoding with the Foundation proof encoding.

\subsection{The precise compilation hypotheses}

Fix a binary one-tape table $M$, including its initial state and distinguished
halt control.  The following are the hypotheses used here.
\begin{enumerate}
\item A finite verified blank-tape startup reaches the prescribed clean
source frame.  Its logical registers are
$(t,w,d,h)=(0,0,0,0)$; the enumeration counter is $2$, and the residual
candidate register is $0$.  The entire finite register file and all spare
cells are part of the frame.
\item Every occurrence of every primitive source instruction has a
whole-tape simulation contract on an arbitrary clean frame.  For a
nonhalting source instruction the contract gives a positive finite physical
run to the correct successor frame; its endpoint control is nonhalting.
For a source halt instruction the contract gives physical halting.  The
contracts cover arbitrary natural register values, gap sizes, and finite
spare tails, not just a bounded sample of them.
\item The source program implements the stated command interpreter and
the complete candidate enumeration.  From a live command boundary each
candidate round takes a positive finite number of source steps, and the
endpoint is the corresponding command result.  A command halt is absorbing.
\item The effective compiler of PA refutations into enumerated command
words is retained or adjusted for the stated axiom schemata, parameter
roles, and head positions.  Every adjustment is a fixed finite proof or
command template, as detailed below.
\end{enumerate}
The source step, source trace predicates, tape step, command updates,
certificate constructors and the finite-sequence codings used in this
argument are primitive recursive.  PA proves the totality of these fixed
functions.  All induction arguments below are carried out on their coded
finite histories using instances of the full PA induction schema.

For the imported PA196 source, the all-model theorem
\texttt{EndpointInternalCodedSourcePhysicalStep.physical\_step} has more
information than its disjunctive headline suggests.  Its proof branches on
whether the resolved source instruction is \texttt{halt}; in the other
branch it constructs positive physical progress and the full successor
frame.  \texttt{PAInternalPhysicalFrame.root\_concrete} shows that every
such boundary has auxiliary control $81$, so it is nonhalting.
\texttt{InternalCodedSourceTrace.trace\_exists} proves totality for source
traces of arbitrary model-internal length.  These are exactly the facts
needed for the reverse argument, rather than an appeal to standard-model
nonhalting.  For a newly compiled table the same facts are supplied by its
own complete symbolic contracts and startup, not by unchanged theorem
names from the old table.

\subsection{Why a closed symbolic runtime certificate is uniform}

The supplied symbolic engines use literal bits, unary runs, and an
arbitrary finite tail of register blocks.  Their soundness has a small
arithmetical proof, which must be understood independently of the
number of tested source sites.
\begin{enumerate}
\item A literal step is one of the fixed table's two rows at its current
state.  When the leading run has affine length $a_0+\sum a_i x_i$,
the checker removes one bit only if $a_0\geq1$ and all coefficients
are nonnegative.  This establishes the side condition for every
natural assignment to the parameters.
\item A unary-run step is used only when the table transition stays
in the same state and moves in the direction of the run.  Induction on
the run length proves that the whole run is traversed and rewritten as
specified, in exactly that many physical steps.  The zero-length case
takes zero steps and leaves the state unchanged.
\item The opaque-list step is justified by a separately rederived
one-block theorem, with arbitrary contexts on both sides and an
arbitrary natural block value.  If a block of value $v$ is traversed
in $v+2$ steps and returns to the same scanning state, induction on
the finite block list gives traversal of the whole list, with duration
equal to its total width.  The empty-list case is the identity.
\item Tape normalization uses literal concatenation, addition of
adjacent run lengths, absorption of a finite zero run into the infinite
zero exterior, and the finite-word identity
$C(A)1=1B(A)$ and its reversal, where
$B(v)=0\,1^{v+1}$ and $C(v)=1\,0\,1^v$.
The last identity follows first for one block and then by induction on
the list $A$.  None of these equalities discards a nonzero spare cell.
\end{enumerate}
Induction on the finite symbolic certificate proves its claimed whole-tape
execution for every natural parameter assignment.  Every successful
primitive contract contains at least one literal physical transition;
unary-run and register-block summaries contribute nonnegative lengths.
Consequently its accumulated physical duration is PA-provably positive.
This uses the checker's moved/progress condition, not merely a positive
count of abstract simplifications.  For a nonhalting instruction the
target is the exact successor frame; a premature halt
cannot be passed over because the table has no transition from halt.
For a halt instruction the certificate ends at halt explicitly.

All four rules are PA-provable facts about coded finite words and lists.
A successful fixed certificate therefore expands into a finite PA proof
of its universally quantified contract.  Bounds on the number of
symbolic inference attempts are limits on finding such a certificate,
not bounds on the natural register values covered by a successful one.
This argument applies to a complete source census and a checked emitter;
it does not upgrade a finite sample of concrete register values into
a universal contract.

\subsection{A bounded reverse-simulation lemma}

\paragraph{Lemma.}
Suppose deterministic physical and source executions start in related
configurations.  Suppose each nonhalting source instruction has a positive
physical simulation ending at a related nonhalting boundary.  If the
physical execution halts in $t$ steps, some one of the first $t+1$ source
instructions is a halt instruction.

\begin{proof}
Assume that the source instructions at indices $0,\ldots,t$ are all
nonhalting.  Induction on $i\leq t+1$ constructs a physical segment
of length $k_i$ ending at the frame for source index $i$, with $k_i\geq i$.
The base length is zero.  The successor contract supplies a length
$\ell_i\geq1$, and concatenation gives $k_{i+1}=k_i+\ell_i\geq i+1$.
The endpoint at $t+1$ is nonhalting and $k_{t+1}\geq t+1$.
By determinism its physical trace agrees, through index $t$, with the
given trace that halts at $t$.  But the longer trace has a next transition
at that index, while the halt control has no outgoing transition.  This
is a contradiction.  Equivalently one may totalize the halt control by
an absorbing self-loop and use halt persistence.

The proof is a PA proof: the induction formula asserts the existence
of coded histories, an endpoint frame, and the inequality $k_i\geq i$.
The defining graphs are arithmetic formulas.  The trace-totality and
contract hypotheses supply the existential witnesses at each successor
stage; full PA induction is applicable to this formula without a bound
on its quantifier complexity.  Trace agreement is proved by induction
up to the shorter trace length.  Nothing in this argument restricts $t$
to an externally standard number.
\end{proof}

With the convention that a source halt instruction sets source control
to zero on its next step, a physical halt by $t$ yields source control
zero by some source index at most $t+1$.  The harmless additive one avoids
conflating ``at a halt instruction'' with ``after executing that instruction.''

The same argument applies once more between candidate boundaries and
source steps.  Write $E_i$ for the command store after $i$ candidate rounds,
starting with candidate $2$.  If a source run reaches control zero by
time $N$, but $E_{N+1}$ is live, then every earlier $E_i$ is live because
the command halt is absorbing.  The positive candidate contracts give a
source run of length at least $N+1$ ending at a live boundary, contradicting
determinism and absorption of source control zero.  Hence some candidate
round, and therefore some finite command history, halts.
For the actual enumeration a round at candidate $n$ uses at most $n+1$
commands.  The first $N+1$ rounds, starting at $2$, therefore use at most
$(N+1)(N+3)$ commands.  This coarse explicit bound suffices for a
primitive recursive extraction: simulate that bounded history with the
proof annotations below and retain the first halting witness.  No
unbounded search for an instruction or command boundary is needed.
The initial four-field command-store encoding is nonzero even though
all four of its fields are zero; it must not be identified with the
distinct halted-command code zero.

The positivity needed at this second level is present in the proof of
the imported \texttt{source\_candidate}: its live case first obtains
\texttt{root\_fresh\_command} with length $m>0$ and appends a consume
trace of length $k\geq0$.  The combined length is $m+k>0$.
\texttt{SourceEnumerationSuffixInduction.source\_from} already supplies
the relevant full PA induction; retaining the additional lower bound
on the accumulated length gives the strengthened assertion used here.

\subsection{The direct interpretation of formulas in PA}

Use the exact constructors
\[
 I(a,b)=4p(a,b),\quad R(a,b)=4p(a,b)+1,\quad
 U(a,B)=4p(B,a)+2,\quad E(a,b)=4p(a,b)+3,
\]
where $p(a,b)=(a+b)(a+b+1)/2+b$.  Code $0$ is truth; code $1=R(0,0)$
is the searched contradiction.  Let $T$ decode these codes directly and
translate $R(a,b)$ to $\operatorname{Bit}(x_a,x_b)$, $E(a,b)$ to
$x_a=x_b$, and $U(a,B)$ to $\forall x_a\,T(B)$.  The first argument
of $\operatorname{Bit}$ is the container.  Code zero translates to $0=0$.
Use a fixed conservative definitional extension of PA for the finitely
many primitive recursive functions needed to express binary digits.
Elimination of these fixed definitions is itself a primitive recursive
proof transformation.

Decoding is primitive recursive by course-of-values recursion, since
every recursive child of a positive formula code is smaller than that
code.  The same induction proves that every variable label in $T(B)$
is strictly below $B$.  Thus $x_B$ is absent from $T(B)$, even bound.
This is a statement about syntax and finite codes, not about truth.

The sole degenerate implication code is $I(0,0)=0$.  Consequently
$T(I(a,b))$ is literally $T(a)\to T(b)$ unless $a=b=0$.
In that exceptional case use a fixed proof of
$(0=0)\leftrightarrow((0=0)\to(0=0))$; an MP step with this code
simply outputs the fixed proof of truth.  PA also proves
\[
 \operatorname{Bit}(u,v)\longrightarrow v<u,
 \qquad \neg\operatorname{Bit}(u,u),
\]
since a set bit implies $2^v\leq u$, whereas $v<2^v$.
Thus the translation of code $1$ has a fixed PA refutation, including
when an enclosing quantifier binds label zero.

\subsection{Uniform PA proofs of every emitted axiom}

For any arithmetic formula $\theta(z,\vec a)$, with the auxiliary
variables fresh, PA proves
\[
 \forall n\,\exists s<2^n\,\forall z<n
 \bigl(\operatorname{Bit}(s,z)\leftrightarrow\theta(z,\vec a)\bigr).
\]
Induct on $n$.  At zero take $s=0$.  At the successor step retain $s$
if $\neg\theta(n,\vec a)$, and take $s+2^n$ if $\theta(n,\vec a)$.
The bound $s<2^n$ ensures that exactly the new digit changes.
The proof is a fixed finite template containing one induction-axiom
instance formed from $\theta$.  Capture-avoiding substitution, fresh-name
selection, construction of this induction instance, and concatenation
of the fixed surrounding proof are primitive recursive.  PA checks
their syntactic correctness uniformly in the code of $\theta$.
No uniform arithmetic truth predicate is being postulated.

For the old gate $G(a,b,c)=\neg(V\to Q)$, let $\varphi=T(c)$ and
assume its guard $c\ne a,b$.  Membership vacuity $V$ has the usual
logical proof.  If $a\ne b$, finite comprehension at cutoff
$x_a+x_b+1$ for
\[
 \bigl(\operatorname{Bit}(x_a,z)\lor z=x_b\bigr)
       \land\varphi[z/x_b]
\]
produces the pointed separation witness.  Renaming the displayed
auxiliary bound variable before substitution prevents capture.
The cutoff contains every eligible point, and the resulting number
has no set bits above the cutoff.  If $a=b$, the first membership
antecedent inside the universally quantified separation formula is
$\operatorname{Bit}(z,z)$ and is false.  In this case use cutoff
$x_b+1$ and the singleton property
$(z=x_b)\land\varphi[z/x_b]$.  This separate branch is necessary:
the occurrence labelled $a$ is also bound when $a=b$.
In either branch the witness proves $\neg Q$, and hence $G$.
If the guard fails, the command's top formula is zero and has the
fixed truth proof.

The six ordinary heads translate to first-order logical or equality
theorems.  Encoded negation and existential quantification use code $1$
as falsity, so their proof templates include the fixed proof of
$\neg\operatorname{Bit}(x_0,x_0)$ rather than treating code $1$ as
an unconstrained predicate.  Coinciding variable labels are handled by
the corresponding finite syntactic branches.  The restricted
propositional head is a tautology; its numeric domination condition
does not affect this soundness direction.

The separate seventh head
\[
 \Gamma(B)=I(B,U(B,B))
\]
has a uniform logical proof because $x_B$ does not occur in $T(B)$.
Its soundness is independent of the gate's guard.  It remains sound
when its source parameter is $A$, $D$, or $H$, and when the seven
heads are emitted in a different fixed order.  Fixed propositional
reassociations or exchanges in the spelling of $G$ have finite
propositional proof templates.

\subsection{A proof-carrying command history, directly in PA}

Let $\operatorname{Prf}_{\rm PA}(q,F)$ be a fixed conventional
primitive recursive proof predicate, with the usual coded constructors
for implication and quantification.  Annotate each live numeric formula
$a$ with a PA proof of $T(a)$.  Annotate the numeric stack by a parallel
stack of such proof codes.  This is ghost data used in the argument;
it adds no tape symbols or control states to $M$.
One may use a proof presentation admitting open theorems.  If the selected
presentation admits only sentences, retain universal closures instead;
the fixed instantiation and closure templates implement the same updates.
Generalization is applied to theorems, not to proofs with undischarged
assumptions.

There is an important zero alias: $p(0,0)=0$, so a pushed truth formula
may disappear into a numerically empty stack.  If
$\operatorname{walk}(w,i)$ is the result of $i$ tail decodings, use
the invariant below, writing $w_i=\operatorname{walk}(w,i)$ and
$q_i=\operatorname{walk}(q,i)$:
\[
 \forall i<w\;
 \bigl[w_i\ne0\ \longrightarrow
 \operatorname{Prf}_{\rm PA}
    (\operatorname{head}(q_i),T(\operatorname{head}(w_i)))\bigr].
\]
For a positive Cantor stack the decoded tail is strictly smaller than
the stack code; hence walking $w$ times reaches zero.  The displayed
bounded invariant therefore controls the entire relevant stack.
After numeric underflow the ghost tail is ignored.  A pop from $w=0$
returns the fixed truth proof and may leave the ghost stack unchanged.
For $w>0$, it pops both stacks.  A push pairs the new formula and proof
with the respective stacks.  These operations preserve the invariant,
including a push of zero onto zero.

Initially the two live formulas and the stack are zero, so the invariant
holds.  AX supplies the six head proofs, the restricted propositional
proof, the generalization of the old head, the saved old top, and the
diagonal proof, in exactly the verified emitter order.  It leaves the
proof of the gate, or truth when the guard fails, in the new top slot;
the new head slot is truth.  MP on its success branch combines the
current two proof codes, using the code-zero exception above if needed.
It assigns the fixed truth proof to the new head and pushes numeric zero
and its fixed truth certificate onto the numeric and ghost stacks,
respectively.  On failure the new top and head are truth and the old
stack is preserved.  Shift moves the proof annotations together with
the numeric fields and uses the conditional pop just described.

All these operations are finite compositions of primitive recursive
code constructors.  Iteration along a coded command history is primitive
recursive, and PA proves preservation by induction on its length.
When Shift halts, $t+h=1$, so PA proves $t=1$ or $h=1$.
The corresponding annotation is a PA proof of
$T(1)=\operatorname{Bit}(x_0,x_0)$.  Append its fixed negation to obtain
a PA refutation.  Thus PA proves that a finite command halt from the
specified initial store yields a PA refutation.

This version needs no converse theorem for
\texttt{Checker.verifyValue}, no conversion between its fuel/tree code
and the HFS table implementation, and no assertion of semantic
soundness of PA.  If the chosen formal library uses a different
conventional PA calculus, a standard primitive recursive proof
translation between the fixed calculi is applied at the end; it is not
a comparison of their truth predicates.

\subsection{The forward direction and permitted source rearrangements}

For the diagonal calculus, the prior proof reconstructs the three
interfaces used by the PA-refutation compiler: bidirectional membership
rewriting, neighborhood equality, and an equality witness.  It uses
$\Gamma(\ulcorner\neg C\urcorner)$ when discharging the separator
for a reverse-rewriting conclusion $C$, and
$\Gamma(\ulcorner E(a,b)\urcorner)$ for the neighborhood interface.
The old gate continues to provide membership vacuity and pointed
separation.  The numerical choices, case distinctions and proof
concatenations are primitive recursive; PA verifies these finite
proof templates.  This is the prior diagonal-calibration argument,
not an inference that a shorter source word is automatically complete.

The bundled and separate-head presentations are intertranslatable:
$\neg(\Gamma(c)\to(V\to Q))$ gives both $\Gamma(c)$ and $G$,
and conversely those two formulas give the bundle.  Every $\Gamma(B)$
is available in the bundled version by taking $c=B$, $a=b=B+1$.
For the separate emitter, a reset of $d,h$ to zero permits arbitrary
values of the three parameters $D,H,A$.  Thus choosing the extra
head as $\Gamma(A)$, $\Gamma(H)$, or $\Gamma(D)$ supplies all diagonal
instances.  Each independent permutation of $D,H,A$ within one of
the same six axiom schemata likewise preserves its full range of
instances.  A permutation of the seven head positions only changes
fixed stack-selection ranks.

There are ten pushed entries: saved top, generalization, restricted
propositional head, and seven ordinary or diagonal heads.
If the seven heads are pushed in order $(h_1,\ldots,h_7)$,
$h_j$ has rank $8-j$ from the stack top, while the propositional head,
generalization and saved top have ranks $8,9,10$.  Starting immediately
after AX, $r+1$ Shifts place rank $r$ in the live top slot.
The prior ten-item compiler identities are
\[
\begin{aligned}
 \mathrm{pop}&=\mathsf{AX}(0,0,0)\,\mathsf S^{11},\\
 \mathrm{push}&=\mathsf{AX}(0,0,0)\,\mathsf S^{9}\,\mathsf{MP}(0),\\
 \mathrm{discard}&=\mathsf{MP}(1)\,\mathsf{AX}(1,0,0)\,
                 \mathsf S^{10}\,\mathsf{MP}(0).
\end{aligned}
\]
They are interpreted with the usual early-halt alternative: encountering
a refutation while preparing a later command already achieves the goal.
All newly inserted heads are implication-coded or negation-coded and
therefore have residue zero modulo four; none is literal false code $1$.
The generalization selection uses ten Shifts.  The ranks of the seven
heads are adjusted to the verified fixed permutation; the other three
ranks do not change.

A single global equality-orientation bit is also harmless.  Define
$E_0(a,b)=4p(a,b)+3$ and $E_1(a,b)=4p(b,a)+3$, and use the chosen
$E_\varepsilon$ consistently in all three equality-using heads and
in the proof templates.  The reverse-rewriting and neighborhood
constructions are schematic in that constructor, and its numerical
code is greater than both labels in either orientation.  The old gate
and the diagonal schema have no separately distinguished equality atom.
Their freshness and proof-supply arguments are unchanged.  In the
arithmetic interpretation the decoded equality is ordinary arithmetic
equality with the corresponding argument order.

The word compiler is then complete for the revised calculus.  Its
reset/compiled-proof/final-Shift word occurs in the complete candidate
enumeration starting at $2$, so PA proves that a PA refutation yields
a halt of the correctly compiled table.  Changing a parameter role
or an axiom's syntax beyond the operations just listed requires a
separate proof transport.  In particular, arbitrary independent
reversals of the primitive equality atom are not justified by assuming
primitive equality symmetry, which the native calculus does not take
as an axiom.

\paragraph{The resulting theorem.}
For any literal table satisfying the four compilation hypotheses and
implementing one of these verified source arrangements, PA proves
\[
 \operatorname{Halt}(M,\mathrm{blank})
   \longleftrightarrow
 \exists p\,\operatorname{Prf}_{\rm PA}(p,\ulcorner\bot\urcorner).
\]
The implication from left to right is the two bounded reverse simulations
followed by the direct PA-proof invariant.  The implication from right
to left is the calibrated refutation compiler and candidate enumeration.
Classical negation gives
\[
 \boxed{\quad
 \mathrm{PA}\vdash
 \neg\operatorname{Halt}(M,\mathrm{blank})
       \longleftrightarrow\operatorname{Con}(\mathrm{PA}).\quad}
\]
This theorem is an ordinary mathematical statement about PA provability.
Applying it to a claimed smaller table requires that table's exact
finite source, all uniform runtime contracts, complete startup, and
forward compiler adjustments to have been verified.  A state-count
candidate lacking any of those ingredients is not covered by the theorem.
